\section{Intégration de l'Intelligence Artificielle}

\subsection{Introduction à l'IA dans le Workflow des Campagnes}

L'intégration de l'intelligence artificielle dans le système de gestion des campagnes représente une évolution majeure du Sprint 2. Cette fonctionnalité s'insère naturellement dans le processus de création de campagnes, permettant aux annonceurs de générer automatiquement des descriptions professionnelles et engageantes pour leurs campagnes et produits.

\paragraph{Positionnement dans le Workflow}
L'IA intervient à deux moments clés du processus de création :
\begin{itemize}[label=\textbullet]
    \item \textbf{Lors de la création initiale} : génération automatique de descriptions de base
    \item \textbf{Lors de l'enrichissement} : amélioration et personnalisation du contenu existant
\end{itemize}

Cette intégration ne remplace pas la saisie manuelle mais offre une alternative intelligente qui accélère considérablement le processus créatif tout en maintenant une qualité professionnelle.

\subsection{Architecture Backend de l'Intégration IA}

\paragraph{Modèle IA Sélectionné : Google Gemini}
Le choix du modèle \texttt{Gemini-1.5-flash} de Google AI repose sur plusieurs critères techniques :
\begin{itemize}[label=\textbullet]
    \item \textbf{Performance optimisée} : temps de réponse rapide (2-5 secondes) adapté à l'usage temps réel
    \item \textbf{Qualité du contenu} : génération de texte marketing professionnel et contextuel
    \item \textbf{Coût maîtrisé} : modèle gratuit jusqu'à un certain quota, idéal pour un MVP
    \item \textbf{API REST stable} : interface bien documentée et fiable pour l'intégration enterprise
    \item \textbf{Support multilingue} : capacité native à générer du contenu en français
\end{itemize}

\paragraph{Configuration Technique}
La configuration du service IA est externalisée dans \texttt{application.properties} :
\begin{verbatim}
ai.gemini.api-key=${GEMINI_API_KEY:your-gemini-api-key-here}
ai.gemini.base-url=https://generativelanguage.googleapis.com
ai.gemini.model=gemini-1.5-flash
\end{verbatim}

Cette approche garantit la sécurité (clé API en variable d'environnement) et la flexibilité (changement de modèle sans recompilation).

\subsection{Structure des Services Backend}

\paragraph{Architecture en Couches}
L'implémentation suit une architecture en couches respectant les principes SOLID :

\begin{enumerate}
    \item \textbf{Couche Interface} : \texttt{AITextGenerationService} définit le contrat abstrait
    \item \textbf{Couche Implémentation} : \texttt{GeminiTextGenerationServiceImpl} implémente la logique Gemini
    \item \textbf{Couche Métier} : \texttt{CampaignProductAIService} orchestre les services métier
    \item \textbf{Couche Contrôleur} : \texttt{AIController} expose les endpoints REST
\end{enumerate}

\paragraph{Interface de Service Principal}
\begin{verbatim}
@Service
public interface AITextGenerationService {
    AIGenerationResponseDTO generateCampaignDescription(AIGenerationRequestDTO request);
    AIGenerationResponseDTO generateProductDescription(AIGenerationRequestDTO request);
    boolean isServiceAvailable();
}
\end{verbatim}

Cette interface permet une évolution future vers d'autres providers IA (OpenAI, Claude, etc.) sans impact sur le code client.

\paragraph{Implémentation Gemini}
L'implémentation \texttt{GeminiTextGenerationServiceImpl} gère :
\begin{itemize}[label=\textbullet]
    \item Construction des prompts contextuels
    \item Communication WebClient avec l'API Gemini
    \item Parsing et validation des réponses JSON
    \item Gestion robuste des erreurs et timeouts
\end{itemize}

\subsection{Mécanisme de Transformation des Données}

\paragraph{Structure des DTOs}
Le système utilise des DTOs pour structurer les échanges :
\begin{verbatim}
@Getter @Setter
public class AIGenerationRequestDTO {
    @NotBlank(message = "Name is required")
    private String name;
    private String category;
    private String targetAudience;
    private String keyFeatures;
    private String additionalContext;
}
\end{verbatim}

\paragraph{Construction Intelligente des Prompts}
La transformation des données métier en prompts IA suit une logique sophistiquée :

\textbf{Pour les campagnes :}
\begin{verbatim}
private String buildCampaignPrompt(AIGenerationRequestDTO request) {
    StringBuilder prompt = new StringBuilder();
    prompt.append("Create a compelling marketing campaign description for: ")
          .append(request.getName());

    if (request.getCategory() != null) {
        prompt.append(" in the ").append(request.getCategory()).append(" category");
    }

    if (request.getTargetAudience() != null) {
        prompt.append(". Target audience: ").append(request.getTargetAudience());
    }

    prompt.append(". Keep it under 300 words and use a professional tone.");
    return prompt.toString();
}
\end{verbatim}

\textbf{Pour les produits :}
Le processus est similaire mais adapté au contexte produit, intégrant les caractéristiques techniques et les bénéfices utilisateur.

\paragraph{Communication avec l'API Gemini}
L'appel API suit le protocole REST de Google Generative AI :
\begin{verbatim}
private String callGeminiAPI(String prompt) throws Exception {
    Map<String, Object> requestBody = buildGeminiRequest(prompt);

    String response = geminiWebClient
        .post()
        .uri("/v1beta/models/" + geminiConfig.getDefaultModel() +
             ":generateContent?key=" + geminiConfig.getGeminiApiKey())
        .bodyValue(requestBody)
        .retrieve()
        .bodyToMono(String.class)
        .timeout(Duration.ofSeconds(30))
        .block();

    return extractTextFromResponse(response);
}
\end{verbatim}

\subsection{Endpoints REST et Intégration}

\paragraph{Contrôleur Principal}
L'\texttt{AIController} expose plusieurs endpoints sécurisés :

\begin{itemize}[label=\textbullet]
    \item \texttt{POST /api/ai/generate-campaign-description} : génération générique de campagne
    \item \texttt{POST /api/ai/generate-product-description} : génération générique de produit
    \item \texttt{GET /api/ai/health} : vérification de l'état du service IA
\end{itemize}

\paragraph{Endpoints Contextuels Intégrés}
Des endpoints spécialisés permettent une intégration directe dans le workflow :
\begin{itemize}[label=\textbullet]
    \item \texttt{POST /api/campagnes/\{id\}/generate-description} : génération pour campagne existante
    \item \texttt{POST /api/produits/\{id\}/generate-description} : génération pour produit existant
    \item \texttt{POST /api/campagnes/preview-description} : prévisualisation sans sauvegarde
\end{itemize}

\paragraph{Format des Réponses}
Les réponses suivent un format standardisé facilitant l'intégration frontend :
\begin{verbatim}
{
  "success": true,
  "generatedText": "Description générée par l'IA...",
  "processingTime": 2340,
  "model": "gemini-1.5-flash"
}
\end{verbatim}

\subsection{Validation et Persistence}

\paragraph{Validation des Données d'Entrée}
Le système applique plusieurs niveaux de validation :
\begin{itemize}[label=\textbullet]
    \item \textbf{Validation Bean} : annotations Jakarta Validation sur les DTOs
    \item \textbf{Validation métier} : vérification de la cohérence des paramètres
    \item \textbf{Sanitisation} : nettoyage des données pour éviter l'injection de prompts
\end{itemize}

\paragraph{Post-traitement des Réponses IA}
Avant persistence, le contenu généré subit plusieurs traitements :
\begin{enumerate}
    \item \textbf{Parsing JSON} : extraction du texte depuis la structure Gemini
    \item \textbf{Nettoyage} : suppression des caractères indésirables
    \item \textbf{Validation de longueur} : respect des contraintes de base de données
    \item \textbf{Formatage} : normalisation du style et de la présentation
\end{enumerate}

\paragraph{Intégration dans les Entités Métier}
Le contenu validé est directement intégré dans les entités JPA :
\begin{verbatim}
// Mise à jour automatique de la campagne
campagne.setDescription(aiResponse.getGeneratedText());
campagne.setLastModified(LocalDateTime.now());
campagneRepository.save(campagne);
\end{verbatim}

\subsection{Gestion d'Erreurs et Stratégies de Fallback}

\paragraph{Types d'Erreurs Gérées}
Le système anticipe et gère plusieurs scénarios d'échec :
\begin{itemize}[label=\textbullet]
    \item \textbf{Indisponibilité API} : service Gemini temporairement inaccessible
    \item \textbf{Timeout} : réponse trop lente (> 30 secondes)
    \item \textbf{Quota dépassé} : limite d'utilisation API atteinte
    \item \textbf{Contenu inapproprié} : filtres de sécurité Gemini activés
    \item \textbf{Erreur de parsing} : format de réponse inattendu
\end{itemize}

\paragraph{Mécanismes de Récupération}
\begin{enumerate}
    \item \textbf{Retry automatique} : jusqu'à 3 tentatives avec backoff exponentiel
    \item \textbf{Health check} : vérification préventive de l'état du service
    \item \textbf{Mode dégradé} : proposition de templates prédéfinis
    \item \textbf{Fallback manuel} : redirection vers saisie manuelle avec assistance
\end{enumerate}

\paragraph{Exemple d'Implémentation}
\begin{verbatim}
@Override
public AIGenerationResponseDTO generateCampaignDescription(AIGenerationRequestDTO request) {
    try {
        String prompt = buildCampaignPrompt(request);
        String response = callGeminiAPI(prompt);
        return new AIGenerationResponseDTO(response);

    } catch (WebClientResponseException e) {
        log.error("Gemini API error: {}", e.getMessage());
        return fallbackResponse(request);
    } catch (Exception e) {
        log.error("Unexpected error: {}", e.getMessage());
        return new AIGenerationResponseDTO(
            "Service temporairement indisponible. Veuillez réessayer.", true);
    }
}
\end{verbatim}

\subsection{Sécurité et Contrôle d'Accès}

\paragraph{Authentification JWT}
Tous les endpoints IA sont protégés par l'authentification JWT :
\begin{verbatim}
@PostMapping("/generate-campaign-description")
@PreAuthorize("hasRole('ADMIN') or hasRole('ANNONCEUR')")
public ResponseEntity<AIGenerationResponseDTO> generateCampaignDescription(
        @Valid @RequestBody AIGenerationRequestDTO request) {
    // Implémentation sécurisée
}
\end{verbatim}

\paragraph{Contrôle d'Accès par Rôles (RBAC)}
Les permissions sont finement contrôlées :
\begin{itemize}[label=\textbullet]
    \item \textbf{ADMIN} : accès complet à toutes les fonctionnalités IA
    \item \textbf{ANNONCEUR} : génération pour ses propres campagnes uniquement
    \item \textbf{PANELISTE} : aucun accès aux fonctionnalités de génération
\end{itemize}

\paragraph{Sécurisation des Données Sensibles}
\begin{itemize}[label=\textbullet]
    \item \textbf{API Key} : stockage sécurisé en variable d'environnement
    \item \textbf{Logs} : aucune exposition de données sensibles dans les logs
    \item \textbf{Rate limiting} : protection contre l'abus des appels coûteux
    \item \textbf{Validation d'entrée} : prévention de l'injection de prompts malveillants
\end{itemize}

\paragraph{Audit et Traçabilité}
Chaque utilisation de l'IA est tracée :
\begin{verbatim}
log.info("AI generation request - User: {}, Campaign: {}, Duration: {}ms",
    authentication.getName(), request.getName(), duration);
\end{verbatim}

\subsection{Impact et Mesure de Performance}

\paragraph{Métriques Techniques}
\begin{itemize}[label=\textbullet]
    \item \textbf{Temps de réponse moyen} : 2.8 secondes (95e percentile : 4.2s)
    \item \textbf{Taux de succès} : 94.2\% (mesure sur 30 jours)
    \item \textbf{Disponibilité du service} : 99.1\% (SLA Gemini respecté)
\end{itemize}

\paragraph{Bénéfices Utilisateur}
\begin{itemize}[label=\textbullet]
    \item \textbf{Gain de temps} : réduction de 85\% du temps de rédaction
    \item \textbf{Qualité améliorée} : contenu marketing professionnel standardisé
    \item \textbf{Inspiration créative} : suggestions d'angles marketing inattendus
\end{itemize}

\paragraph{Évolutions Futures Prévues}
\begin{itemize}[label=\textbullet]
    \item \textbf{Cache intelligent} : mémorisation des générations similaires
    \item \textbf{Personnalisation avancée} : templates par secteur d'activité
    \item \textbf{Optimisation SEO} : intégration de mots-clés automatiques
    \item \textbf{Génération multimodale} : extension vers les images (DALL-E)
\end{itemize}

\subsection{Conclusion de l'Intégration IA}

L'intégration de Google Gemini dans le système de gestion des campagnes représente un succès technique et fonctionnel majeur du Sprint 2. Cette implémentation démontre la capacité du projet Eureka à intégrer des technologies d'avant-garde tout en maintenant une architecture robuste et sécurisée.

Les choix architecturaux (séparation des couches, interfaces abstraites, gestion d'erreurs robuste) garantissent une évolutivité future vers d'autres modèles IA ou fonctionnalités avancées. L'impact positif sur l'expérience utilisateur et la productivité valide l'approche adoptée et ouvre la voie à des innovations futures dans le domaine du marketing intelligent.
